\begin{abstract}
Investors with assets in several currencies often first choose the asset portfolio and then manage its currency risk
with a currency overlay, a portfolio of currency positions. Choosing the asset and currency positions together in one
step can never do worse in sample, but it may fit noise in the data. This thesis studies whether such joint
optimization improves on the sequential currency overlay out of sample. We compare the two approaches under five risk
measures, the variance, semi-variance, mean absolute deviation, Expected Shortfall, and second-order Expected Shortfall.
We backtest them in three investable universes of increasing size, with and without regularization, and with a mean
term added to the unregularized programs. We find that, without a mean term, the joint and overlay optimizers achieve
very similar risk-adjusted returns. The in-sample advantage of joint optimization in the minimized risk measure is realized out of sample only in the smallest universe.
For Expected Shortfall and second-order Expected Shortfall, this advantage comes from fitting the currency positions to
the worst days in the data and from an implicit reward for a high realized mean, and neither generalizes. With a mean
term, the joint optimizers achieve higher Sharpe ratios than the overlay optimizers for every risk measure in the two
larger universes. All optimizers have much lower volatility and smaller drawdowns than the equally weighted portfolio,
and higher risk-adjusted returns in the two larger universes. Regularization and the mean term improve their
risk-adjusted returns further. Managing the currencies actively lowers risk relative to full hedging at similar
risk-adjusted returns and outperforms leaving them unhedged in the two larger universes. These results suggest that
investors who only minimize risk can keep their currency overlay unless their universe is small, while those who also
use expected returns can benefit from joint optimization.

    % Multi-currency portfolio construction requires managing asset and currency exposures together, yet it 
    % remains unclear whether these decisions are best made simultaneously or sequentially as the investable universe grows. 
    % In this thesis, we build on the work of \cite{RegularizedMultiCurrencyExpectedShortfallPortfolios,SparseAndStableInternationalPortfolioOptimization}. 
    % We investigate whether joint optimization of asset and currency weights provides a benefit over the standard sequential currency overlay, 
    % across five risk measures and three universe sizes. The empirical backtests show that joint optimization does not deliver a meaningful 
    % improvement in realized excess returns or risk-adjusted returns over the overlay, in any configuration tested, 
    % with the possible exception of the smallest universe. This motivates a closer 
    % look at where joint optimization differs from the overlay by 
    % construction, which is risk minimization. The thesis identifies two mechanisms behind the joint program's in-sample risk advantage for the tail 
    % measures, a tail-fitting mechanism and an implicit risk-return tradeoff, and shows that neither generalizes out of sample. For the deviation measures, 
    % the in-sample advantage is small and most of it survives out of sample. The out-of-sample advantage of joint optimization is positive only in 
    % the small universe, where the currency weights are a larger share of the decision variables.
\end{abstract}